% Kodierung und Sprache
\usepackage[utf8]{inputenc}
\usepackage[T1]{fontenc}
\usepackage[ngerman]{babel}

% Schrift: Helvetica (wie Arial in Google Docs)
\usepackage{helvet}
\renewcommand{\familydefault}{\sfdefault}

% Seitenformat
\usepackage[a4paper, margin=2.5cm]{geometry}

% Mathematik
\usepackage{amsmath}
\usepackage{amssymb}
\usepackage{mathtools}

% Diagramme
\usepackage{tikz}
\usetikzlibrary{calc, positioning, patterns}

% Farben für Balkendiagramme
\usepackage{xcolor}
\definecolor{parteiA}{gray}{0.25}
\definecolor{parteiB}{gray}{0.65}
\definecolor{parteiC}{RGB}{100,160,220}
\definecolor{parteiD}{RGB}{80,160,80}
\definecolor{parteiE}{RGB}{200,60,60}
\definecolor{parteiF}{RGB}{220,190,40}
\definecolor{parteiS}{RGB}{160,120,40}
\definecolor{diagrammHG}{RGB}{210,230,245}

% Tabellen
\usepackage{booktabs}
\usepackage{array}

% Listen
\usepackage{enumitem}

% Typografie
\usepackage{microtype}
\usepackage{csquotes}

% Hyperlinks
\usepackage[colorlinks=true, linkcolor=blue, urlcolor=blue, citecolor=blue]{hyperref}

% Grafiken (für Screenshots-Platzhalter im Anhang)
\usepackage{graphicx}

% Hilfsmakro: Skalierung für Balkendiagramme (100\% = 8\,cm)
\newlength{\balkenhoehe}
\setlength{\balkenhoehe}{8cm}
\newcommand{\proztocm}[1]{\dimexpr#1\balkenhoehe/100\relax}
